\section{Proofs}\label{section2}

\subsection[Asymptotics of $\Gamma^\C$]{Asymptotics of $\boldsymbol{\Gamma^\C}$}\label{ss:asy}
Let us derive formula \eqref{eq:asy1} for the asymptotics of the function $\Gamma^C(a+\xi|a'-\ov \xi)$.
First, let us verify that the right-hand side is single-valued
on $\Lambda$.
We must verify that the expression
\[
\exp\big\{ \xi\ln\xi - \ov\xi\ln\ov\xi\big\}, \qquad \xi=\tfrac{k+{\rm i}s}2
\]
is single valued.
Represent $\xi=r {\rm e}^{{\rm i}\phi}$. We transform our expression as
\begin{gather*}
\exp\big\{\xi(\ln(r)+{\rm i}\phi+2\pi {\rm i} n)-
\ov \xi(\ln(r)-{\rm i}\phi-2\pi {\rm i} n) \big\} \\
\qquad{}=\exp\big\{(\xi-\ov\xi) \ln r +(\xi+\ov\xi)({\rm i}\phi+2\pi {\rm i} n) \big\}.
\end{gather*}
We have $\xi+\ov\xi\in \Z$ and therefore the exponential is single valued.

Next, we apply the Stirling formula in the form
(see \cite[equation~(1.18(3))]{HTF1})
\[
\ln
\Gamma(c+z)=\big(z+c-\tfrac12\big)\ln z -z+\tfrac 12\ln(2\pi)+O\big(z^{-1}\big),
\qquad \text{where $|\arg (z)|<\pi-\epsilon$},
\]
and get
\begin{gather*}
\Gamma(a+\xi) \simeq \sqrt{2\pi}\exp
\big\{ \big(\xi+a-\tfrac12\big)\ln\xi-\xi \big\},\\
\Gamma(1-a'-\ov \xi) \simeq \sqrt{2\pi}\exp
\big\{ \big(\ov\xi-a'+\tfrac12\big)\ln\ov\xi-\ov\xi\big\}.
\end{gather*}
The ratio gives us \eqref{eq:asy1}.
An evaluation of the asymptotics in the sector $|\arg(-z)|<\pi-\epsilon$
gives the same result.

\begin{Corollary}\label{cor:L2}
Let $\Re(a_\alpha+a_\alpha')>0$, $\Re(b_\beta+b_\beta')>0$.
Denote
\[
\upsilon:=\sum (a_\alpha+a_\alpha')+\sum(b_\beta+b_\beta').
\]
\begin{enumerate}\itemsep=0pt
\item[$(a)$] If $\upsilon<p+q$, then $F(k,{\rm i}s):=
\Kc{p}{q}\big[\genfrac{}{}{0pt}{}{(a|a')}{(b|b')};\tfrac{k+{\rm i}s}{2}\bigl|;\tfrac{-k+{\rm i}s}{2}\big]$
tends to $0$ as $|k+{\rm i}s|$ tends to $\infty$.
\item[$(b)$]
If $\upsilon<p+q-1$, then for each $k$ the function $F(k,{\rm i}s)$ as a
function in $s$ is integrable.
Under the same condition $F(k,{\rm i}s)$ is contained in $L^2(\Lambda)$.
\item[$(c)$]
If $\upsilon<p+q-2$, then $F(k,{\rm i}s)$ is contained in $L^1(\Lambda)$.
\end{enumerate}
\end{Corollary}

\begin{Remark}
Formula \eqref{eq:asy1} also gives us asymptotics of
$\Kc{p}{q}$ on vertical lines $\sigma=h+{\rm i}s$. Indeed,
\[
\Gamma\big(a+\tfrac{k+h+{\rm i}s}{2}\big|a'+\tfrac{-k+h+{\rm i}s}{2}\big)=
\Gamma\big(a+\tfrac h2+\tfrac{k+{\rm i}s}{2}\big|a'+\tfrac h2+\tfrac{-k+{\rm i}s}{2}\big),
\]
and we can control integrability under shifts of the integration contour.
\end{Remark}

\subsection{Decomposition of Mellin--Barnes integrals in residues}\label{ss:meijer}
Now we start to prove Theorem~\ref{th:1}, i.e., to derive
the quadratic expressions of the func\-tions~$\Gc{p}{q}$ in terms of
the usual hypergeometric functions $\hyp{p}{q}$.

First, we write the definition \eqref{eq:def} of $\Gc{p}{q}$ in the form
\[
\sum_{k\in \Z} z^{-k/2\big|k/2} I_k(z), \qquad
I_k(z)=\frac{1}{2\pi {\rm i}}\int (\dots)|z|^{-\sigma}\,{\rm d}\sigma.
\]

The integrals $I_k(z)$ are special cases of
Mellin--Barnes integrals, i.e., integrals of the type
\begin{equation}
J^{A,B}_{C,D}(a,b,c,d;u)=\frac 1{2\pi {\rm i}}\int_L \frac{\prod\limits_{\alpha=1}^A\Gamma(a_\alpha+\sigma)
\prod\limits_{\beta=1}^B\Gamma(b_\beta-\sigma)}
{\prod\limits_{\gamma=1}^C\Gamma(c_\gamma+\sigma)\prod\limits_{\delta=1}^D\Gamma(d_\delta-\sigma)} u^{-s}
\, {\rm d}s,
\label{eq:J}
\end{equation}
where $L$ is a contour separating poles of factors
$\Gamma(a_\alpha+\sigma)$ and poles of $\Gamma(b_\beta-\sigma)$.
The behavior of such Mellin--Barnes integrals
(Meijer $G$-functions) was investigated by Meijer,
his results are exposed in \cite{BW, Luke,Mar}.
Under certain conditions $J(a,b,c,d;u)$ can be expressed in terms
of functions $\Sigma_\pm(z)$, where
$\Sigma_+(z)$ is the sum of residues of the integrand at poles of
factors $\Gamma(a_\alpha+\sigma)$, and $\Sigma_-(z)$
the sum of residues at poles of $\Gamma(b_\beta-\sigma)$.

In our case $A=D=q$, $B=C=p$, and $u=|z|$ is a positive real.
The contour $L$ coincides with the imaginary axis at infinity.
The asymptotics of the absolute value of the integrand on the imaginary axis
is
\[
\sim |\sigma|^{\sum\limits_{\alpha=1}^q\Re (a_\alpha-d_\alpha)+\sum\limits_{\beta=1}^p\Re(b_\beta- c_\beta)},
\]
see \cite[equation~(1.18(4))]{HTF1}.
One of the statements of \cite[Theorem~18]{Mar} implies that
\emph{if the integrand tends to zero on the imaginary axis,
then the integral $J(a,b,c,d;u)$ conditionally converges if
the integrand tends to $0$ at infinity for $u>0$ and is given by}
\begin{itemize}\itemsep=0pt
\item[--] \textit{$\Sigma_+(u)$ for $q>p$;}
\item[--] \textit{$\Sigma_-(u)$ for $q<p$;}
\item[--] \textit{for $p=q$ we have $\Sigma_+(u)$ if $0<u<1$ and $\Sigma_-(u)$ for $u>1$.}
\end{itemize}
If
$\sum\Re(a_\alpha-d_\alpha)+\sum\Re(b_\beta-c_\beta)<-1$,
then the absolute convergence is obvious.

In our case $a_\alpha$ is replaced by $a_\alpha+k/2$, $b_\beta$ by
$b_\beta-k/2$, and $c_\beta\mapsto 1-b'_\beta-k/2$,
$d_\alpha\mapsto 1-a'_\alpha+k/2$.
Therefore under our conditions \eqref{eq:cond1}, \eqref{eq:cond2}
all integrals $I_k$ converge.

\subsection{Evaluation of sums of residues}\label{ss:evaluation}

\begin{Lemma}
Fix $a|a'$ and consider the following family of functions:
\[
\gamma_k(\sigma):=\Gamma^\C\big(a+\tfrac {k+\sigma}2\big|a'+\tfrac
{-k+\sigma}2\big).
\]
Denote by $\Omega$ the sets of all points $(k,\sigma)\in \Z\times \C$
such that $\sigma$ is a pole of $\gamma_k(\sigma)$.
\begin{enumerate}\itemsep=0pt
\item[$(a)$] The set $\Omega$ is contained in the half-plane $\Re\sigma<0$ if and only
if $\Re(a+a')>0$.
\item[$(b)$] The set $\Omega$ consists of points
\begin{gather*}
k=-m+m'-a+a', \\
\sigma= -m-m'-a-a',
\end{gather*}
where $m$, $m'$ range in $\Z_+$.
\end{enumerate}
\end{Lemma}

\begin{proof} (a)
Let all poles be contained in the domain $\Re\sigma<0$.
Taking $k=a'-a$ we get
\[
\gamma_{a'-a}(\sigma)=\Gamma^\C
\big(a+\tfrac {a'-a+\sigma}2\big|a'+\tfrac {-a'+a+\sigma}2\big)=
\Gamma^\C\big(\tfrac{a+a'}2 +\tfrac{\sigma}{2}\big|\tfrac{a+a'}2
+\tfrac{\sigma}{2}\big).
\]
The point $\sigma= -(a+a')$ is a pole of $\gamma_{a'-a}(\sigma)$.
Therefore, $\Re(a+a')>0$.

Conversely, let $\Re(a+a')>0$.
Consider $(k,\sigma)\in \Omega$.
Then (see Section~\ref{ss:Gamma})
\begin{equation}
a+\tfrac{k+\sigma}2=-m\in \Z_-, \qquad a'+\tfrac{-k+\sigma}2=-m'\in \Z_-.
\label{eq:Omega}
\end{equation}
Therefore $a+a'+\sigma\in \Z_-$ and $\Im\sigma<0$.

(b) We solve the system \eqref{eq:Omega}.
\end{proof}

\begin{proof}[Evaluation of the sum of residues (the proof of Theorem~\ref{th:1})]
For definiteness assume that $p\le q$.
Let us write the residue $R(j;m,m')$ of the integrand \eqref{eq:def}
at a point
\[
k= -m+m'-a_j+a'_j,\qquad \sigma=-m-m'-a_j-a'_j.
\]
We have
\[
\tfrac{k+\sigma}{2}=-a_j-m,\qquad \tfrac{-k+\sigma}{2}=-a_j'-m',
\]
and
keeping in mind \eqref{eq:res-gamma} we get
\begin{gather*}
R(j;m,m')=\frac{2(-1)^m}{m!\,m'!}
 \prod_{\alpha\ne j}
\Gamma^\C\big(a_\alpha-a_j-m\big|a_\alpha'-a_j'-m'\big)\\
\phantom{R(j;m,m')=}{} \times
\prod_{\beta} \Gamma^\C\big(b_\beta+a_j+m\big|b_\beta'+a_j'+m'\big)
\times z^{a_j+m|a_j'+m'}.
\end{gather*}
Applying \eqref{eq:shift1} and \eqref{eq:shift2} we come to
\begin{gather*}
R(j;m,m')=2 \frac{(-1)^m}{m!\,m'!}
\prod_{\alpha\ne j} \frac{\Gamma^\C\big(
a_\alpha-a_j\big|a_\alpha'-a_j'\big)(-1)^{m}}
{(1-a_\alpha+a_j)_m(1-a_\alpha'+a_j')_{m'}} \\
\hphantom{R(j;m,m')=}{} \times
\prod_{\beta} \Gamma^\C\big(b_\beta+a_j\big|b_\beta'+a_j'\big)
b_\beta+a_j\big)_m(b_\beta'+a_j')_{m'} (-1)^{m'}
\times z^{a_j+m|a_j'+m'}.
\end{gather*}
Reordering factors, we get
\begin{gather*}
2 z^{a_j|a'_j}
\prod_{\alpha\ne j} \Gamma^\C\big(a_\alpha-a_j\big|a_\alpha'-a_j'\big)
\prod_{\beta}
\Gamma^\C\big(b_\beta+a_j\big|b_\beta'+a_j'\big) \\
\qquad{}\times
(-1)^{mq} \frac{z^m \prod\limits_{\beta}(b_\beta+a_j\big)_{m}}
{m!\prod\limits_{\alpha\ne j} (1-a_\alpha+a_j)_{m}}
 \times
(-1)^{m'p} \frac{\ov z^{m'} \prod\limits_{\beta}(b_\beta'+a_j'\big)_{m'}}
{m'!\prod\limits_{\alpha\ne j} (1-a_\alpha'+a_j')_{m'}}.
\end{gather*}

Formal calculation with series gives
\begin{gather*}
\sum_{m,m'} R(j,m,m')=
2 z^{a_j|a_j'}
 \prod_{\alpha\ne j} \Gamma^\C\big(a_\alpha-a_j\big|a_\alpha'-a_j'\big)
 \prod_{\beta}
 \Gamma^\C\big(b_\beta+a_j\big|b_\beta'+a_j'\big)
\\
\hphantom{\sum_{m,m'} R(j,m,m')=}{} \times
 \sum_m \frac{ \big((-1)^{q} z\big)^m \prod\limits_{\beta}(b_\beta+a_j\big)_m} {m!\prod\limits_{\alpha\ne j} (1-a_\alpha+a_j)_m}
 \times
 \sum_{m'} \frac{ \big((-1)^{p} \ov z\big)^{m'} \prod\limits_{\beta}(b_\beta+a_j\big)_{m'}}
 {m'!\prod\limits_{\alpha\ne j} (1-a_\alpha'+a_j')_{m'}} .
 \end{gather*}
We get a product of two hypergeometric series whose radius of convergence
is $\infty$ if $q>p$ and $1$ if $q=p$.
They are absolutely convergent in the disk of convergence,
therefore the last identity really takes place.
Therefore
\[
\sum_j \sum_{m,m'} R(j,m,m')=\Sigma^\C_+ (z).
\]
On the other hand the absolute convergence allows us to write
 \begin{equation}
 \sum_{m,m'} R(j,m,m')=\sum_{k\in\Z}\,\,
\sum_{m,m'\colon m-m'=k} R(j,m,m')=\sum_{k\in\Z} (z/\ov z)^{k/2} I_k.
 \label{eq:sum-R}
 \end{equation}
Thus we proved the coincidence of $\Gc{p}{q}$ and $\Sigma_+^\C$.
The summation of residues at the right poles of
$_p^{\vphantom{\C}}\cK^\C_q$ is similar.

It is important that the convergence of the series $\sum\limits_{k\in\Z}$
in~\eqref{eq:sum-R} is locally uniform in the parameters
$a_\alpha$, $b_\beta$ near any point $(a_0)\in\C^q$, $(b_0)\in\C^p$,
for which the coefficients at $z^l\ov z^{l'}$ are well-defined.
\end{proof}
